\section{Introduction}

\subsection{Contexte}

Dans un code de calcul de structures par éléments finis (code\_aster, Abaqus, FEniCS, etc.), le comportement du matériau est calculé \emph{localement}, en chaque point d'intégration (point de Gauss) de chaque élément. À chaque itération de l'équilibre global, le code fournit à la loi de comportement la déformation en fin de pas et l'état des variables internes au début du pas. La loi lui renvoie deux choses :
\begin{enumerate}[nosep]
  \item la contrainte et les nouvelles variables internes ;
  \item l'\emph{opérateur tangent} $\partial\sig_{n+1}/\partial\eps_{n+1}$, que le Newton global utilise pour assembler sa matrice de rigidité.
\end{enumerate}

Écrire cette tangente à la main est la partie la plus longue et la plus sujette à erreur du développement d'une loi. Une tangente fausse ne fausse pas la solution (le résidu reste exact), mais elle dégrade la convergence : un calcul qui devrait converger en 3 itérations en demande 30, ou diverge. La \emph{différentiation automatique} (AD) supprime ce travail. Les frameworks d'apprentissage automatique comme JAX apportent en plus la compilation (\code{jit}) et la vectorisation (\code{vmap}) sur CPU et GPU.

Ce mini-projet reproduit, en réduction, la chaîne complète : écrire une loi en JAX, valider sa tangente, l'utiliser dans un Newton, l'intégrer par lots et mesurer ses performances. Il se termine par un bonus : un potentiel hyperélastique représenté par un réseau de neurones et identifié sur des données expérimentales.

\subsection{Contributions}

\begin{itemize}[nosep]
  \item Une loi J2 en notation de Mandel avec écrouissage isotrope linéaire ou de Voce et cinématique de Prager. Elle est intégrée par retour radial, avec un Newton local sur $\Delta p$ dont la dérivée est obtenue au choix à travers les itérations ou par le théorème des fonctions implicites (\code{jax.custom\_jvp}).
  \item La démonstration et la correction de trois pièges : \code{if} Python, NaN dans les gradients en mode inverse (« double \code{where} »), perte des symétries mineures en notation $3\times3$.
  \item Trois tangentes (AD, analytique, différences finies) dont l'accord est vérifié par 77 tests \code{pytest}.
  \item Un pilote de point matériel en chargement mixte. Il montre que l'écrouissage linéaire en chargement proportionnel ne permet pas d'observer la convergence quadratique, et explique pourquoi.
  \item Un benchmark NumPy / JAX CPU / JAX GPU (calcul seul et avec transferts) en FP64 et FP32, analysé avec un modèle roofline et mesuré sur une RTX~2050 portable, avec un kit de mesure réutilisable (dossier \code{gpu/}).
  \item Le calcul de $\Pk = \partial\psi/\partial\F$ et $\mathbb{A}=\partial^2\psi/\partial\F^2$ par \code{grad}/\code{hessian}, puis l'identification d'un potentiel neuronal polyconvexe sur les données de Treloar.
\end{itemize}

\subsection{Environnement de mesure}

Deux machines ont servi. Les sections 7 à 10 et les mesures CPU de la section~\ref{sec:perf} ont été faites sur une machine virtuelle à 2 cœurs Intel Xeon (2,1~GHz) et 7~Go de mémoire, sans GPU, avec Python~3.13. Les mesures CPU y sont donc \emph{relatives} (rapports NumPy/JAX, fraction de la bande passante atteinte). Les mesures GPU (section~\ref{sec:gpu}) ont été faites sur un portable Lenovo LOQ (i5-12450HX, RTX~2050 4~Go, Ubuntu~24.04, JAX~0.12 avec CUDA~13). Les deux machines utilisent JAX en double précision.

\subsection{Organisation du rapport}

Les sections~\ref{sec:maths} et~\ref{sec:meca} rassemblent les rappels mathématiques et mécaniques nécessaires. Elles peuvent être sautées par un lecteur familier du sujet. La section~\ref{sec:integration} dérive l'algorithme de retour radial et la section~\ref{sec:tangente} la tangente cohérente. La section~\ref{sec:jax} présente la différentiation automatique, JAX et ses pièges. Les sections~\ref{sec:verification} à~\ref{sec:perf} présentent les résultats : vérification, point matériel, performances. La section~\ref{sec:hyper} traite le bonus hyperélastique.
